% Exact isolated candidate proof for source-aware verification.
% This file makes no novelty or priority claim and is not a live-book input.

\begin{proof}
\medskip\noindent\emph{Imported source scope.}
In \cite[Example~6.7.2]{CLD25}, Chambert--Loir and Ducros put a
\(G\)-tropical affine structure on an interval
\(I\subseteq\overline{\mathbf R}\), with
\(A=\partial I\cap\overline{\mathbf R}\), by requiring parlinear functions
to be constant near every point of \(A\).  In particular, every affine
function on a connected domain meeting \(A\) is constant.  Fix a finite
endpoint \(a\in A\cap\mathbf R\).  Together with their definition of smooth
forms in \cite[Definition~6.2.1]{CLD25}, the example says the following.
By \cite[\S\S6.1.1 and 6.1.6]{CLD25}, a tropical chart is affine and each of
its scalar coordinate functions is affine.  After shrinking to a connected
neighbourhood meeting \(A\), Example~6.7.2 therefore makes every coordinate
constant.  The chart, and every smooth \(0\)-form pulled back through it, is
constant there.  Use the ordinary inward real coordinate on the finite
stratum and extend it topologically by the value \(t(a)=0\); this is not a
nonconstant affine chart through \(a\).  In this coordinate, the smooth
functions are precisely the functions described by
\eqref{eq:p071-endpoint-collar-smoothness}.  Conversely, use the constant
terminal chart on the collar and the ordinary coordinate \(t\) on
endpoint-free interior neighbourhoods; these local descriptions glue.  On
the one-dimensional interior,
\[
 \mathrm d'\mathrm d'' f
   =f''(t)\,\mathrm d't\wedge\mathrm d''t,
\]
so positivity is exactly \eqref{eq:p071-endpoint-stage-zero}; the form
vanishes on the collar where \(f\) is constant.  Definition~8.1.1 and
Lemma~8.1.4 of \cite{CLD25} give the generated sheaf and its ordinal
filtration used in the statement.  We will also use the elementary
uniform-to-monotone device of \cite[Lemma~8.1.10]{CLD25}: if real-valued
functions \(f_j\) converge uniformly to \(f\) on a set \(K\), then
\[
                  f_j+2\lVert f_j-f\rVert_{K}
\]
has a decreasing subsequence converging uniformly to \(f\).

\medskip\noindent\emph{The comparison class is a sheaf.}
Restrictions in \(\mathscr J_{\partial}\) are immediate.  We check
gluing.  Suppose that \(W=\bigcup_{\nu}W_\nu\) and that
\(u|_{W_\nu}\in\mathscr J_{\partial}(W_\nu)\) for every \(\nu\).  Let
\(C\) be a connected component of \(W\).  The function \(u|_C\) is locally
extended-real convex.

We first show that \(u|_C\) is either finite everywhere or identically
\(-\infty\).  Let
\[
 Z=\{t\in C:u(t)=-\infty\},
 \qquad
 F=C\setminus Z.
\]
If \(t\in Z\) is an interior point of \(C\), choose an interval
\((t-\delta,t+\delta)\cap C\) on which \(u\) is convex.  For a point
\(s>t\) sufficiently close to \(t\), choose \(q>s\) in that interval and
write \(s\) as a strict convex combination of \(t\) and \(q\).
Then \eqref{eq:p071-endpoint-extended-convexity} gives \(u(s)=-\infty\).
The same argument on the left shows that \(Z\) is open at \(t\).  If
\(0\in Z\), choose \(a>0\) such that \(u\) is convex on \([0,a)\).
For \(0<s<a\), choose \(q\in(s,a)\) and write
\(s=(1-s/q)0+(s/q)q\); again \(u(s)=-\infty\).  Thus \(Z\) is open also
at the endpoint.

If \(t\in F\) is an interior point and there were points \(s\in Z\)
arbitrarily close to \(t\), choose such an \(s\) and a point \(q\) on the
other side of \(t\), all in one interval of local convexity.  Expressing
\(t\) as a strict convex combination of \(s\) and \(q\) would give
\(u(t)=-\infty\), a contradiction.  Hence \(F\) is open at every interior
point.  If \(0\in F\), the local instance of
\eqref{eq:p071-endpoint-boundary-minimum} excludes the value
\(-\infty\) on a whole initial interval, so \(F\) is open at \(0\) as
well.  Since \(C\) is connected, either \(Z=C\) or \(F=C\).

In the finite case, the standard one-dimensional local-to-global
criterion shows that \(u|_C\) is convex.  Here is a direct verification.
Given \(a<c<b\) in \(C\), cover \([a,b]\) by finitely many intervals on
which \(u\) is convex.  Take a Lebesgue number \(\delta>0\) for this finite
cover of \([a,b]\), and choose a partition
\[
             a=x_0<x_1<\cdots<x_N=b
\]
that contains \(c\) and satisfies
\(x_{i+1}-x_{i-1}<\delta\) whenever the three terms occur.  Each interval
\([x_{i-1},x_{i+1}]\) then lies in one member of the cover.  Local convexity
therefore says that the successive secant slopes are nondecreasing.  Summing
the resulting inequalities on the two sides of \(c\) yields
\[
 \frac{u(c)-u(a)}{c-a}
 \leq
 \frac{u(b)-u(c)}{b-c},
\]
which is equivalent to the convexity inequality at \(c\).  Thus
\eqref{eq:p071-endpoint-extended-convexity} holds on every interval in
\(C\).  It is automatic when \(u|_C\equiv-\infty\).

It remains to propagate the boundary minimum.  Suppose \(0\in C\), fix
\(t\in C\), and choose \(s\in(0,1)\) so small that \(st\) lies with \(0\)
in one member of the cover.  There
\[
                              u(0)\leq u(st).
\]
If \(u(0)=-\infty\), the desired inequality is automatic.  If
\(u(0)\in\mathbf R\), the preceding dichotomy makes \(u\) finite on
\(C\), while global convexity gives
\[
                 u(st)\leq(1-s)u(0)+s u(t).
\]
Combining the last two inequalities, cancelling \((1-s)u(0)\), and
dividing by \(s\) gives \(u(0)\leq u(t)\).  This proves the gluing axiom
and hence that \(\mathscr J_{\partial}\) is a sheaf.

\medskip\noindent\emph{The generated sheaf is contained in the comparison
class.}
The sheaf \(\mathscr J_{\partial}\) is stable under pointwise decreasing
limits.  Indeed, every inequality
\eqref{eq:p071-endpoint-extended-convexity} passes to the limit in
\(\mathbf R\cup\{-\infty\}\), since no \(+\infty\) value occurs, and
\eqref{eq:p071-endpoint-boundary-minimum} passes to the limit as well.

It also contains \(\operatorname{Psh}_{S,0}\).  On an interior interval,
\eqref{eq:p071-endpoint-stage-zero} gives ordinary convexity.  On the
component meeting \(0\), the derivative of a generator is zero on the
constant collar and is nondecreasing thereafter; consequently the
generator is convex on the whole half-interval and takes its minimum at
\(0\).  By the minimal definition of the generated psh sheaf,
\begin{equation}
                       \operatorname{Psh}_S
                       \subseteq\mathscr J_{\partial}.
 \label{eq:p071-endpoint-generated-in-jensen}
\end{equation}

\medskip\noindent\emph{The extended-real alternatives at the endpoint.}
Let \(0\in W\) and \(u\in\mathscr J_{\partial}(W)\).  After shrinking its
connected component, choose \(R>0\) with \([0,R)\subseteq W\).

Suppose first that \(u(0)=-\infty\).  For \(0<t<R\), choose
\(q\in(t,R)\).  Since
\[
                         t=(1-t/q)0+(t/q)q,
\]
extended-real convexity gives
\[
 u(t)\leq(1-t/q)u(0)+(t/q)u(q)=-\infty.
\]
Thus \(u\equiv-\infty\) on \([0,R)\).  The smooth psh constants \(-n\)
decrease to this function on the fixed neighbourhood.

Assume now that
\begin{equation}
                              u(0)\in\mathbf R.
 \label{eq:p071-endpoint-finite-value}
\end{equation}
The boundary minimum excludes \(-\infty\) at every other point, so \(u\)
is real-valued.  For \(0<s<t<R\), convexity and
\eqref{eq:p071-endpoint-boundary-minimum} give
\begin{equation}
 u(s)
 \leq\left(1-\frac{s}{t}\right)u(0)+\frac{s}{t}u(t)
 \leq u(t).
 \label{eq:p071-endpoint-monotonicity}
\end{equation}
Hence \(u\) is nondecreasing.  It is continuous in the interior by
ordinary finite convexity.  At the endpoint, for any fixed \(t\in(0,R)\),
\eqref{eq:p071-endpoint-boundary-minimum} and the first inequality in
\eqref{eq:p071-endpoint-monotonicity} show that
\[
 u(0)\leq u(s)
 \leq u(0)+\frac{s}{t}\bigl(u(t)-u(0)\bigr)
 \quad(0<s<t),
\]
so \(u(s)\to u(0)\) as \(s\downarrow0\).  Thus \(u\) is continuous on
\([0,R)\).

\medskip\noindent\emph{A fixed-domain collar regularization.}
Choose
\[
                         0<r<R_1<R_2<R.
\]
Put
\[
                         m=
             \frac{u(R_2)-u(R_1)}{R_2-R_1}\geq0
\]
and define a finite function \(\bar u:\mathbf R\to\mathbf R\) by
\begin{equation}
 \bar u(t)=
 \begin{cases}
  u(0),&t\leq0,\\
  u(t),&0\leq t\leq R_1,\\
  u(R_1)+m(t-R_1),&t\geq R_1.
 \end{cases}
 \label{eq:p071-endpoint-convex-extension}
\end{equation}
This extension is convex and nondecreasing.  At \(0\), the left slope is
zero and every right secant slope is nonnegative by
\eqref{eq:p071-endpoint-boundary-minimum}.  At \(R_1\), convexity of
\(u\) shows that every left secant slope is at most the secant slope
\(m\) to \(R_2\).  The one-dimensional slope criterion for convexity
therefore applies at both joining points.

Fix a nonnegative even function
\(\theta\in C^\infty_c((-1,1))\) with integral one, and put
\(\theta_\sigma(s)=\sigma^{-1}\theta(s/\sigma)\).  For a sequence
\(\sigma_j\downarrow0\), define
\begin{equation}
                         A_j=\bar u*\theta_{\sigma_j}.
 \label{eq:p071-endpoint-convolution}
\end{equation}
Every \(A_j\) is smooth, convex, and nondecreasing.  Evenness gives
\(\int s\theta_{\sigma_j}(s)\,\mathrm ds=0\), so Jensen's inequality
gives
\begin{equation}
                              A_j\geq\bar u.
 \label{eq:p071-endpoint-convolution-dominates}
\end{equation}
Moreover, \(A_j\to\bar u\) uniformly on every compact interval.

We next impose exact constancy on a collar.  For \(\eta>0\), set
\begin{equation}
 M_\eta(a,b)=
 \int_{\mathbf R^2}
 \max\{a+s_1,b+s_2\}
 \theta_\eta(s_1)\theta_\eta(s_2)\,
 \mathrm ds_1\,\mathrm ds_2.
 \label{eq:p071-endpoint-regularized-maximum}
\end{equation}
This function is smooth, convex, and nondecreasing in each variable.
The zero first moments and the support of the kernel give
\begin{equation}
 \max\{a,b\}
 \leq M_\eta(a,b)
 \leq\max\{a,b\}+\eta.
 \label{eq:p071-endpoint-regularized-max-bounds}
\end{equation}
If \(a>b+2\eta\), then \(a+s_1>b+s_2\) throughout the convolution
support and \(M_\eta(a,b)=a\); symmetrically,
\begin{equation}
 M_\eta(a,b)=b
 \quad\text{whenever}\quad b>a+2\eta.
 \label{eq:p071-endpoint-regularized-max-truncation}
\end{equation}

Choose \(\varepsilon_j\downarrow0\), put
\(\eta_j=\varepsilon_j/4\), and define
\begin{equation}
 H_j(t)=M_{\eta_j}(t,\varepsilon_j),
 \qquad
 F_j(t)=A_j\bigl(H_j(t)\bigr).
 \label{eq:p071-endpoint-collar-approximants}
\end{equation}
The function \(H_j\) is smooth, convex, and nondecreasing.  From
\eqref{eq:p071-endpoint-regularized-max-bounds} one gets, for \(t\geq0\),
\begin{equation}
 0\leq H_j(t)-t\leq\varepsilon_j+\eta_j.
 \label{eq:p071-endpoint-reparametrization-error}
\end{equation}
By \eqref{eq:p071-endpoint-regularized-max-truncation},
\begin{equation}
 H_j(t)=\varepsilon_j
 \quad\text{for}\quad
 0\leq t<\varepsilon_j-2\eta_j
                =\frac{\varepsilon_j}{2}.
 \label{eq:p071-endpoint-exact-collar}
\end{equation}
Consequently \(F_j\) is smooth and exactly constant on an initial
interval.

The function \(F_j\) is convex.  Indeed, \(A_j'\geq0\),
\(A_j''\geq0\), \(H_j'\geq0\), and \(H_j''\geq0\), so
\[
 F_j''
  =A_j''(H_j)(H_j')^2+A_j'(H_j)H_j''
  \geq0.
\]
Thus \(F_j|_{[0,r)}\) belongs to
\(\operatorname{Psh}_{S,0}([0,r))\).  It also dominates \(u\) on
\([0,r]\): by monotonicity of \(A_j\),
\eqref{eq:p071-endpoint-reparametrization-error}, and
\eqref{eq:p071-endpoint-convolution-dominates},
\begin{equation}
 F_j(t)=A_j(H_j(t))
 \geq A_j(t)
 \geq\bar u(t)=u(t).
 \label{eq:p071-endpoint-smooth-domination}
\end{equation}

Finally, \(F_j\to u\) uniformly on \(K=[0,r]\).  For all sufficiently
large \(j\), both \(t\) and \(H_j(t)\) lie in the fixed compact interval
\([0,r+1]\).  Hence
\begin{align}
 \sup_{t\in K}|F_j(t)-u(t)|
 &\leq
 \sup_{x\in[0,r+1]}|A_j(x)-\bar u(x)|
 \notag\\
 &\quad+
 \sup_{t\in K}|\bar u(H_j(t))-\bar u(t)|
 \longrightarrow0,
 \label{eq:p071-endpoint-uniform-convergence}
\end{align}
where the second term tends to zero by uniform continuity of \(\bar u\)
and \eqref{eq:p071-endpoint-reparametrization-error}.

Let
\[
                       e_j=\lVert F_j-u\rVert_K.
\]
By \cite[Lemma~8.1.10]{CLD25}, the sequence \(F_j+2e_j\) has a
subsequence, say
\[
                         G_n=F_{j_n}+2e_{j_n},
\]
which decreases uniformly to \(u\) on \(K\).  Adding a constant preserves
smoothness, collar constancy, and convexity.  Restricting to the one fixed
open neighbourhood \(V=[0,r)\) proves
\eqref{eq:p071-endpoint-fixed-domain-approximation} in the finite case.

\medskip\noindent\emph{Interior points and conclusion.}
Let \(x\in W\cap(0,R_0)\).  On a sufficiently small connected interval
about \(x\), the dichotomy proved above says that a section of
\(\mathscr J_{\partial}\) is either identically \(-\infty\) or is a
finite convex function.  The first case is the decreasing limit of the
smooth psh constants \(-n\).  In the second case, choose a compact
interval \(K=[a,b]\) around \(x\), contained in a slightly larger interval
on which \(u\) is finite and convex.  The one-sided derivatives of a finite
convex function are finite at the interior points \(a,b\).  Choose
\[
                    p_-\in\partial u(a),
                    \qquad p_+\in\partial u(b),
\]
and extend \(u|_{[a,b]}\) by
\[
 u(a)+p_-(t-a)\quad(t<a),
 \qquad
 u(b)+p_+(t-b)\quad(t>b).
\]
Monotonicity of the subgradients makes this a finite convex function on
\(\mathbf R\).  Convolution with the even kernels used in
\eqref{eq:p071-endpoint-convolution} gives smooth convex functions
converging uniformly to \(u\) on \(K\).  Applying
\cite[Lemma~8.1.10]{CLD25} and restricting to the interior of \(K\)
produces a decreasing sequence of smooth convex functions on one fixed
neighbourhood of \(x\).

We have proved
\[
                 \mathscr J_{\partial}
                 \subseteq\operatorname{Psh}_{S,1}.
\]
The tautological inclusion
\(\operatorname{Psh}_{S,1}\subseteq\operatorname{Psh}_S\), together with
\eqref{eq:p071-endpoint-generated-in-jensen}, proves
\eqref{eq:p071-endpoint-first-stage-collapse}.  The construction at the
endpoint gives every assertion following
\eqref{eq:p071-endpoint-fixed-domain-approximation}.
\end{proof}
